\section{Description of the Problem}
\label{sec:problem}

Section~\ref{sec:introduction} established that building semantic models unlock significant value
but are costly to create.
This section examines the cost problem in concrete terms.
We describe what BACnet points are, quantify the labor required to map them, and identify the research
gap that has prevented an automated solution.

\subsection{The BACnet Point Mapping Problem}
\label{subsec:bacnet-problem}

A Building Automation System (BAS) exposes operational data through a hierarchy of \emph{points}, 
individual sensors, actuators, setpoints, alarms, and status flags.
In BACnet (the dominant open protocol for commercial building controls), each point is a typed object
(Analog Inputs, Binary Outputs, Analog Values, etc.)\ with a unique address on the network.
A typical commercial office building has between 500 and 5{,}000 BACnet points; a large hospital or
campus facility may have tens of thousands~\cite{Trenbath2022}.

The \emph{mapping problem} is this: a raw list of BACnet points contains no standard information about
what physical equipment each point belongs to, what role it plays in the system, or how that equipment
connects to other components.
A point named AHU-1.SA-T tells a human engineer ``probably supply air temperature on
Air Handling Unit 1,'' but only if the engineer already knows the naming convention used by that
particular BMS installer.
It is common to find conventions that differ between contractors, between projects, and even between
systems within the same building.

Reconstructing the full picture of the building, which points belong to which equipment, what each
point means, and how the equipment connects into a complete HVAC system, requires reading as-built
HVAC drawings, interrogating BMS programming, and applying expert engineering judgment.
This process is manual, time-consuming, and dependent on domain expertise.

Modbus, an even older and simpler protocol used in legacy HVAC controllers, meters, and industrial
equipment, exposes data as bare numbered registers with no descriptive metadata at all, making the
mapping challenge even more difficult than BACnet.

\subsection{The Cost of Manual Mapping}
\label{subsec:cost}

Manual BACnet point mapping and BAS commissioning are widely documented as major cost drivers in
building construction and retrofit projects.
Table~\ref{tab:commissioning-costs} summarizes the key figures from industry research and published
case studies.

\begin{table}[htbp]
\centering
\caption{Key cost and effort data for manual BAS commissioning and point mapping.}
\label{tab:commissioning-costs}
\small
\begin{tabular}{p{5.5cm}p{3.5cm}p{4cm}}
\toprule
\textbf{Cost Item} & \textbf{Figures} & \textbf{Source} \\
\midrule
BAS commissioning as \% of construction & 0.5\%--1.5\% of total; 8\%--15\% of BAS project value
  & \cite{BCA2024} \\
\addlinespace
Case study: 150{,}000\,sq\,ft school
  & \$225{,}000 commissioning (\$1.50/sq\,ft)
  & \cite{PingCx2024} \\
\addlinespace
Complete BAS installation average
  & \$7.76/sq\,ft (range: \$5.20--\$14.18)
  & \cite{Trenbath2022} \\
\addlinespace
Legacy BMS integration premium
  & +20\%--40\% over base cost
  & \cite{Envigilance2023} \\
\addlinespace
Additional service calls from poor mapping
  & 12--15 extra calls/year; +30\%--45\% troubleshooting time
  & \cite{PingCx2024} \\
\addlinespace
Energy performance drift from poor commissioning
  & 15\%--30\% energy waste in first 3 years; \$0.20--\$0.50/sq\,ft/year
  & \cite{PingCx2024} \\
\addlinespace
First-year unplanned expenses (case study school)
  & \$175{,}000, nearly equal to original commissioning cost
  & \cite{PingCx2024} \\
\addlinespace
HVAC technician shortage (2024--2025)
  & 110{,}000 unfilled positions
  & \cite{AccessCoins2025} \\
\bottomrule
\end{tabular}
\end{table}

To illustrate the compounding nature of these costs: the 150{,}000\,sq\,ft school in the case study
paid \$225{,}000 for commissioning up front, then incurred approximately \$175{,}000 in unplanned
first-year expenses due to inadequate point documentation, nearly doubling the effective cost of the
commissioning phase~\cite{PingCx2024}.
According to NREL~\cite{Trenbath2022}, complete BAS installation averages \$7.76 per square foot
across the United States, with the legacy systems integration premium adding a further 20 to 40
percent on retrofit projects~\cite{Envigilance2023} (data from 2022 and 2023).

Manual approaches are also incomplete by design.
Industry practice typically tests only 10 to 20 percent of terminal units during commissioning~\cite{PingCx2024};
every 10 percent reduction in testing coverage correlates to approximately a 6 to 8 percent increase
in post-occupancy issues.
The result is a pervasive, chronic problem: buildings begin their operational life with
under-documented, partially verified control systems, and the cost of correction grows over time. 
Wang et al.~\cite{Wang2017} describe manual point mapping as ``labor intensive and costly, presenting
a major impediment to innovations in building control'', validating the scale of the problem even
where precise figures are contractor-dependent.

\subsection{The Research Gap and Emerging Solutions}
\label{subsec:research-gap}

Despite the scale and economic impact of the manual mapping problem, no widely-deployed end-to-end
automated solution existed as of Wang et al.'s 2017 systematic review~\cite{Wang2017}, which found
the field lacked:
\begin{itemize}[leftmargin=*, itemsep=2pt]
  \item a complete problem formulation with integrated solution approaches,
  \item standardized test datasets and performance metrics, and
  \item methods well-aligned with real-world mapping requirements.
\end{itemize}

The AI wave of 2023 onwards has catalyzed a new generation of attempts to close this gap.
Table~\ref{tab:landscape} surveys the most relevant commercial and academic efforts and positions
SI-Mapper within the landscape.
The Building Commissioning Association's published guidance~\cite{BCA2024} confirms that
point-to-point verification remains a mandated step on every new BAS project, underscoring that the
problem persists even as partial solutions emerge.

\begin{table}[htbp]
\centering
\caption{Competitive landscape for automated BAS point mapping and semantic model generation (2024--2026).}
\label{tab:landscape}
\small
\begin{tabular}{p{3cm}p{3cm}p{3cm}p{4cm}}
\toprule
\textbf{Solution} & \textbf{Category} & \textbf{Input} & \textbf{Gap vs.\ SI-Mapper} \\
\midrule

PingCx~\cite{PingCx2024}
  & Automated functional commissioning \& point-to-point testing
  & Live BACnet connection
  & Tests whether points respond correctly; does \emph{not} build a semantic model or identify equipment topology \\
\addlinespace

BrickLLM~\cite{Perini2025}
  & AI-driven Brick ontology generation
  & Free-text building description
  & Requires accurate natural-language input; does not parse BACnet CSV exports, HVAC drawings, or produce 223P \\
\addlinespace

BuildingMOTIF~\cite{BuildingMOTIF_NREL}
  & DOE open-source SDK for 223P / Brick / Haystack model creation and validation
  & Expert-authored RDF; BACnet point lists (manual)
  & Developer toolkit, not an AI inference engine; requires domain expert to drive model creation \\
\addlinespace

BIM2RDF / Open223~\cite{Open223}
  & BIM-to-223P conversion
  & Revit / IFC BIM file
  & Requires a complete, accurate BIM as input; does not handle buildings where only as-built drawings and BACnet CSVs exist \\
\addlinespace

Willow~\cite{Willow2025}
  & AI digital twin platform with ML device classification
  & Live BMS / IoT feeds
  & Uses proprietary DTDL ontology, not 223P; requires live data connection and human-guided onboarding; \$110M-funded commercial platform \\
\addlinespace

SkySpark / Project Haystack~\cite{SkySpark}
  & Analytics platform with semantic tagging
  & Engineer-applied Haystack tags
  & Tagging is manual; provides analytics engine once tagged, not an automated mapping tool \\
\addlinespace

\textbf{SI-Mapper (this work)}
  & \textbf{End-to-end AI pipeline for 223P model generation}
  & \textbf{BACnet CSV point exports + as-built HVAC drawings}
  & \textbf{No comparable system} addresses the full pipeline from raw BACnet exports and engineering drawings to a validated ASHRAE 223P graph \\

\bottomrule
\end{tabular}
\end{table}

The table reveals a clear structural gap.
Existing solutions either stop short of semantic modeling (PingCx automates testing but produces no
ontology), require clean inputs that do not exist in practice (BrickLLM needs accurate text;
BIM2RDF needs a complete BIM file), or are expert-facing development toolkits rather than turnkey
automation (BuildingMOTIF).
Commercial platforms such as Willow approach the problem from the opposite direction, top-down,
starting from live data in large portfolios, using a proprietary ontology that is not aligned with
ASHRAE 223P or the emerging interoperability standards required for Virtual Power Plants integration.

National laboratory programs (NREL, LBNL, PNNL, NIST) have made significant progress in defining
the 223P standard and associated tooling through BuildingMOTIF~\cite{BuildingMOTIF_NREL,DOE_BuildingMOTIF},
but the explicit DOE framing acknowledges the goal is still to ``lower costs and reduce installation
time'', an admission that the current tooling does not yet achieve this without substantial expert
involvement.

The gap SI-Mapper addresses is therefore specific and narrow:
\emph{automated end-to-end generation of ASHRAE 223P semantic models from the artifacts that
actually exist on the ground}, BACnet CSV exports and as-built engineering drawings, without
requiring a pre-existing BIM file, clean natural-language descriptions, or an expert semanticist.

\subsection{Why This Matters for the Energy Grid}
\label{subsec:grid-relevance}

The point mapping problem is not only a building operations challenge.
It is a direct barrier to the energy transition.

\textbf{Non-Wire Alternatives} and \textbf{Virtual Power Plants} aggregate distributed
energy resources, including commercial building HVAC loads, to provide grid services equivalent to
building new power plants or transmission lines.
HVAC systems consume approximately 40 percent of total commercial building energy and can be
modulated to respond to grid signals, shifting load away from peak demand periods.
In aggregate, this flexibility represents enormous potential grid capacity.

The Virtual Power Plants market is growing rapidly.
North American Virtual Power Plants programs totaled 37.5\,GW of flexible behind-the-meter capacity
as of 2025, up approximately 14 percent from 2024, with active deployments rising more than 33
percent in the same period~\cite{WoodMac2025}.
Policy support is accelerating: ten state legislatures introduced Virtual Power Plants bills in
2024~\cite{RMI_VP3_2025,SEPA_NCCETC2025}; Massachusetts utilities are implementing a
\$50\,million Grid Services Compensation Fund to promote non-wire alternatives~\cite{RMI_VP3_2025};
and the California Energy Commission awarded a \$2\,million EPRI grant specifically to demonstrate
Virtual Power Plants for commercial buildings including HVAC controls~\cite{CEC_EPRI2024}.

\textbf{How aggregators operate today.}
Two types of aggregators participate in these programs.
\emph{Hardware aggregators} manufacture controllable devices (smart thermostats, batteries, EV
chargers) and aggregate their own installed base into a dispatchable resource portfolio.
\emph{Software aggregators} act as protocol proxies, connecting heterogeneous devices from
multiple vendors under a single management layer.
Both types sell capacity and demand response services to utilities and grid operators, who dispatch
them using standardized signals such as OpenADR or IEEE~2030.5.
This model works efficiently for residential assets and fleets of identical devices because the
aggregator controls the endpoint hardware and already knows its capabilities.

\textbf{Commercial buildings are structurally different.}
A commercial building's HVAC system is not a device the aggregator manufactured.
It was designed by an engineering firm, installed by a mechanical contractor, and programmed by a
BMS provider (Siemens, Johnson Controls, ABB, Niagara, or equivalent).
Each of these parties has a legitimate interest in protecting the reliability of the system: the BMS
provider guarantees uptime; the commissioning engineer designed the control sequences; the building
operator is legally responsible for occupant comfort.
None of them will willingly allow an external party to write arbitrary setpoints to their system
without a formal agreement and a clear technical boundary.
The result is that every commercial building enrollment today requires a custom, building-by-building
negotiation and integration project \cite{DOE_VPP2025}.
The DOE Virtual Power Plants Liftoff report identifies that no industry standardization
exists. Virtual Power Plants vendors have been building their solutions from the ground up, then
working on a case-by-case basis with utilities to develop customized solutions, resulting in rigid
systems that can't easily be replicated in new areas~\cite{DOE_VPP2025}.

\textbf{How an ASHRAE 223P ontology changes the logistics.}
An ASHRAE 223P semantic model of a building is not a control system replacement, it is a
machine-readable contract that describes what is controllable, how to reach it, and under what
conditions.
Consider the operational chain it enables:

\begin{enumerate}[leftmargin=*, itemsep=2pt]
  \item \textbf{Asset discovery.} The ontology declares every controllable component (fans, coils,
        VAV boxes, dampers) as a typed entity with its position in the HVAC topology.
        An aggregator can query it to answer ``what loads does this building contain and what are
        their nominal capacities?'' without a site visit or manual documentation review.
  \item \textbf{Live telemetry access.} ASHRAE 223P's \emph{external reference} mechanism links
        each semantic property directly to its BACnet object address.
        A supply air temperature sensor in the model carries the address of the physical BACnet
        point that reports its value.
        The aggregator reads live sensor data by following those references, with no custom
        integration code per building.
  \item \textbf{Flexibility estimation.} With topology and live telemetry available, an 
        algorithm can compute how much load the building can shed, for how long, and with what
        thermal lag, before comfort constraints are violated.
        This estimate feeds directly into the aggregator's dispatch optimization.
  \item \textbf{Controlled dispatch.} The ontology also identifies the BACnet addresses of
        controllable setpoints (supply air temperature setpoint, fan speed command, cooling
        setpoint).
        When the grid operator sends an OpenADR curtailment signal, the aggregator translates it
        into targeted BACnet writes at precisely the addresses the ontology specifies, within
        the boundaries agreed contractually with the building owner.
        The BMS continues to run its reliability logic; the aggregator adjusts only the setpoints
        that were explicitly licensed under the agreement.
\end{enumerate}

